\documentclass[11pt]{article}

\usepackage[margin=1in]{geometry}
\usepackage{booktabs}
\usepackage{array}
\usepackage{amsmath}
\usepackage{hyperref}

\title{RelaLeap: Early Evidence for Contextual Support Routing in Sparse Columnar Residual Adapters}
\author{Draft generated from the RelaLeap automation evidence trail}
\date{2026-06-24}

\begin{document}
\maketitle

\begin{abstract}
RelaLeap studies small, auditable residual adapters added to a frozen
sequence model. The working hypothesis is that a sparse set of residual
columns can learn reusable corrections while disturbing the frozen base less
than dense or poorly routed alternatives. The early harness phase succeeded:
experiments are command-driven, reproduce across local and Colab backends, and
emit checked artifacts. The strongest positive scientific result so far is not
an inference-time settling result. It is a support-routing result. Replacing a
linear top-k support router with a contextual MLP support router produced large
cross-entropy reductions on both character-level and tokenized Tiny
Shakespeare settings, expanded support usage, and reduced support churn. In
the bounded promotion gate, contextual routing improved larger-character
alpha-0 cross entropy from 3.2079 to 1.8480 in Colab and tokenized larger
alpha-0 cross entropy from 3.5254 to 2.8681. The effect was replicated locally
and in Colab, and the contextual router was promoted as the default
support-wide router. However, the evidence does not yet justify a strong claim
that individual residual columns are causally separable reusable mechanisms:
rank-matched top-k-1 contextual routing currently has better cross entropy than
the promoted top-k-2 contextual run, and held-out fingerprint stability remains
to be measured. The useful interim conclusion is that support selection is a
major bottleneck, contextual token-neighborhood features are a powerful way to
route sparse residual columns, and causal-column claims require stricter
deconfounding.
\end{abstract}

\section{Motivation}

The goal of RelaLeap is not to produce a large language model benchmark win in
the first campaign. The goal is to build a reproducible methodology for asking
whether residual columns can provide modular, reusable, low-interference
corrections to a frozen sequence model.

The central research question after the first harness phase is:
\[
\text{Can a sparse residual layer learn causally separable reusable
corrections with less interference than matched alternatives?}
\]
Cross entropy and perplexity are retained as guardrails, but are not the only
scientific target. The more interesting quantities are support selection,
oracle-support regret, functional churn, causal intervention fingerprints,
continual-learning retention, and comparisons against dense or random-support
controls.

\section{Implemented System}

\subsection{Frozen Base and Residual Adapter}

The current experiments use small Tiny Shakespeare sequence models. The base is
a frozen two-layer sequence model; hidden dimension and sequence length vary
by experiment scale. A residual adapter is inserted at one site. The adapter is
a sparse set of residual columns. Each token position chooses a top-k subset
of columns, mixes the selected column values, and adds the result to the base
hidden state.

In the implementation, each column contains several learned residual atoms.
The atom values are initialized to zero. This makes the residual adapter an
exact identity map at initialization while still allowing gradients to flow
into residual parameters. For a hidden vector \(h_t\), the adapter computes
column scores \(s_c(h_t)\), chooses a support set \(S_t\), and adds
\[
r(h_t) =
\sum_{c \in S_t}
\operatorname{softmax}_{c \in S_t}(s_c(h_t)) v_c,
\qquad
h'_t = h_t + r(h_t).
\]
Here \(v_c\) is the atom-weighted value vector for column \(c\). Early runs
used \(k=1\). Support-wide runs use \(k=2\).

\subsection{Linear Versus Contextual Support Routing}

The original router was linear:
\[
s(h_t) = W h_t.
\]
The contextual MLP router uses a richer local token-neighborhood feature
vector. The implemented feature vector concatenates the current hidden state,
the previous hidden state, the next hidden state, first differences, and
position features:
\[
\phi_t =
\left[
h_t,\ h_{t-1},\ h_{t+1},\ h_t-h_{t-1},\
h_{t+1}-h_t,\ p_t,\ \sin(2\pi p_t),\ \cos(2\pi p_t)
\right].
\]
Column scores are then produced by a layer-normalized MLP:
\[
s(h_t) = W_2\,\operatorname{GELU}(W_1\,\operatorname{LN}(\phi_t)).
\]
In the main promotion-gate settings the contextual router hidden width was
128, the base hidden dimension was 96, there were 24 residual columns, four
atoms per column, and \(k=2\).

\subsection{HEP Settling}

The code also implements hidden-state error propagation (HEP), an
inference-time settling procedure controlled by an alpha sweep. HEP variants
included pinned support, clipped residual-adapter updates, a guided supervised
oracle, prediction-entropy settling, and temporal-consistency settling.

Temporal-consistency clipped HEP passed a promotion gate as a safe label-free
support-stress mitigation, but the effect size is small. In the larger-char and
tokenized local/Colab promotion gate, the mean loss improvement from nonzero
temporal HEP over alpha 0 was \(8.231\times 10^{-5}\), with very small logit
deltas. In the contextual-router experiments, nonzero HEP alphas did not
explain the large win; alpha 0 was best after contextual routing. Thus the main
positive result is a training-time support-routing result, not a settling-time
result.

\subsection{Artifact Discipline}

Every accepted experiment is command-driven and writes standard artifacts:
\texttt{summary.json}, \texttt{metrics.csv}, and \texttt{notes.md}. After
completion, the same artifact checker validates the result tree. Decision
reports summarize whether a gate is passed and whether a default should be
changed. This mattered because many early HEP effects were extremely small and
because Colab execution was brittle.

\section{Experimental Progress}

\subsection{Phase 0: Harness and Initial HEP}

The first campaign validated a minimal char-level Tiny Shakespeare harness.
The base had two layers, hidden dimension 32, sequence length 32, one
insertion site, eight residual columns, four atoms per column, \(k=1\), seed 1,
and 10 residual training steps. The initial HEP alpha sweep showed that
\(\alpha=0.25\) improved cross entropy while staying within a logit-delta
budget. Larger alphas improved raw loss more but violated the stability
budget.

\begin{center}
\begin{tabular}{rrrr}
\toprule
Alpha & CE loss & Loss improvement vs alpha 0 & Max logit delta \\
\midrule
0.00 & 3.56317067 & 0.00000000 & 0.00000000 \\
0.25 & 3.55195642 & 0.01121426 & 0.05185765 \\
0.50 & 3.54104090 & 0.02212977 & 0.10371491 \\
1.00 & 3.52012682 & 0.04304385 & 0.20742965 \\
\bottomrule
\end{tabular}
\end{center}

This established that HEP could move the model in a useful direction, but also
that unconstrained HEP was too easy to overdrive.

\subsection{Support Stress and Clipped HEP}

The next step intentionally created support repicking under support-stress
conditions. Pinned support exposed large instability: support change fraction
reached 0.7891, and pinned-versus-repicked logit delta reached 3.3260.
Clipped HEP reduced the pinned-versus-repicked logit delta to about 0.00235,
but simple clipping alone did not improve support-stress loss.

A guided clipped HEP oracle then showed that useful hidden-state directions
exist: supervised guidance during settling improved support-stress loss inside
the stability budget. Because it used labels during settling, it was treated as
an oracle, not as a deployable inference mechanism.

Label-free HEP objectives were then tested. Prediction entropy was stable but
poor. Temporal consistency was stable and repeatedly slightly helpful. It was
promoted as the default support-stress mitigation, but its practical effect is
small. This motivated shifting attention from settling objectives to support
selection and column structure.

\subsection{Residual Objective Variants}

After temporal-clipped HEP was promoted as a support-stress mitigation, several
residual training objectives were tested. The default supervised
cross-entropy residual objective remained strongest under the current gates.
PC logit-MSE, anchored PC, confidence penalty, margin penalty, label smoothing,
focal CE, and train-time temporal consistency did not justify a default change.

For example, in the objective-discriminative local/Colab PC gate, supervised CE
had mean best HEP loss 3.58667445, while PC logit-MSE had mean best HEP loss
3.59590077. PC improved its own objective but did not beat supervised residual
training on supervised CE.

\subsection{Support Width Deconfounding}

The first major clue that support selection mattered came from a controlled
support-width matrix. The experiment fixed supervised CE training,
temporal-clipped HEP, disabled the support-stress preset, used sequence length
64 and hidden dimension 64, and crossed column count with top-k:

\begin{center}
\begin{tabular}{lrrrrrr}
\toprule
Variant & Columns & Top-k & Best HEP CE & Used cols & Unique supports & Max frac \\
\midrule
Baseline & 12 & 1 & 3.58667445 & 1 & 1 & 1.0000 \\
Support width & 12 & 2 & 3.49384785 & 10 & 15 & 0.3750 \\
Capacity & 24 & 1 & 3.58667445 & 1 & 1 & 1.0000 \\
Capacity + support & 24 & 2 & 3.49411678 & 15 & 17 & 0.3594 \\
\bottomrule
\end{tabular}
\end{center}

Increasing top-k from 1 to 2 improved best HEP CE by about 0.0928 and expanded
support utilization from one used column to ten. Increasing stored columns
from 12 to 24 while keeping \(k=1\) did nothing. This suggested that active
support width, not merely stored capacity, was the current bottleneck.

\subsection{Exhaustive Support Audits}

At validation scale, all \( \binom{12}{2} = 66 \) two-column supports could be
evaluated exactly. The learned router was compared to the per-token oracle
support and to global fixed supports. The per-token oracle beat the learned
router on most positions:

\begin{center}
\begin{tabular}{lr}
\toprule
Quantity & Value \\
\midrule
Router loss & 3.49828529 \\
Per-token oracle loss & 3.46230602 \\
Oracle-support regret & 0.03597923 \\
Oracle-support positive fraction & 0.87301588 \\
Router minus best global fixed support loss & -0.03413963 \\
Contextual holdout oracle-gap recovery & 0.75938903 \\
Strongest pairwise synergy & 0.06710577 \\
\bottomrule
\end{tabular}
\end{center}

The router still beat the best single global fixed pair, which pointed toward
better token-conditioned routing rather than replacing routing with one fixed
support. Hidden-only selectors did poorly, but contextual selectors recovered a
large fraction of the oracle gap.

The same pattern repeated on larger 24-column support-width audits. On
larger-char seed 2, the contextual support intervention recovered 0.8930 of
the router-oracle holdout gap and reduced holdout CE by 0.0861 versus the
learned router. On seed 3, it recovered 0.9056 of the gap and reduced holdout
CE by 0.1950. A train-time-style contextual support head also helped, with gap
recoveries 0.8612 and 0.7693 on seeds 2 and 3.

\section{Main Positive Result: Contextual MLP Support Routing}

The support audits suggested that the router, not the residual atoms alone, was
the bottleneck. The project therefore implemented and tested a learned
contextual MLP support router. The first local and Colab larger-char
comparisons showed large improvements over the linear top-k-2 router:

\begin{center}
\begin{tabular}{lrrrrrr}
\toprule
Backend & Linear CE & Context CE & Delta & Linear used & Context used & Churn delta \\
\midrule
Local & 3.09682608 & 1.91732013 & -1.17950594 & 12 & 18 & -0.45507812 \\
Colab & 3.15306783 & 1.91732001 & -1.23574781 & 13 & 18 & -0.35156250 \\
\bottomrule
\end{tabular}
\end{center}

This result was strong but initially came from one scale. A bounded promotion
gate then required a larger-char seed-2 repeat and a tokenized larger setting,
with both local and Colab evidence. The matrix preserved supervised CE,
temporal-clipped HEP, \(k=2\), support-stress preset disabled, and the same
artifact checks.

\begin{center}
\begin{tabular}{llrrrrr}
\toprule
Backend & Dataset & Linear CE & Context CE & Delta & Used cols & Unique supports \\
\midrule
Local & char, seed 2 & 3.15261722 & 1.84800947 & -1.30460775 & 19 & 60 \\
Local & word, seed 1 & 3.51685524 & 2.86810708 & -0.64874816 & 20 & 50 \\
Colab & char, seed 2 & 3.20794559 & 1.84802544 & -1.35992014 & 19 & 60 \\
Colab & word, seed 1 & 3.52541637 & 2.86810851 & -0.65730786 & 20 & 52 \\
\bottomrule
\end{tabular}
\end{center}

Every gate cell passed. Contextual routing lowered alpha-0 cross entropy in
all four cells, expanded support utilization in all four cells, and reduced
support churn in all four cells. Nonzero temporal-clipped HEP did not win in
any of the four cells. The default change was therefore scoped to the support
router, not to HEP settling.

\section{Post-Promotion Check}

After the contextual router was promoted as the support-wide default, fresh
local promoted-default artifacts were generated for validation-char,
larger-char, and tokenized larger settings:

\begin{center}
\begin{tabular}{lrrrr}
\toprule
Experiment & Columns & Used cols & Unique supports & Final CE \\
\midrule
Validation char & 12 & 11 & 30 & 3.09003305 \\
Larger char & 24 & 18 & 44 & 1.91732013 \\
Tokenized larger & 24 & 20 & 50 & 2.86810708 \\
\bottomrule
\end{tabular}
\end{center}

Alpha 0 remained best in all three runs. Again, this supports the interpretation
that the useful mechanism is train-time contextual support routing, not
test-time HEP improvement.

\section{Causal Column Fingerprints: Promising but Not Yet Proven}

The latest audit asked whether the promoted contextual top-k-2 residual layer
has interpretable causal column fingerprints. The audit measured loss changes
under column ablation and forced-column interventions, then bracketed the
result with controls.

\begin{center}
\begin{tabular}{lrrrr}
\toprule
Variant & CE & Residual norm & Used cols & Unique supports \\
\midrule
Learned top-k-2 contextual & 2.91240144 & 4.43444967 & 19 & 41 \\
Rank-matched top-k-1 contextual & 2.84867883 & 4.67442513 & 33 & 33 \\
Random fixed top-k-2 & 3.81906438 & 4.16884756 & 2 & 1 \\
Dense rank/FLOP norm-matched & 3.25567937 & 4.43444967 & -- & -- \\
\bottomrule
\end{tabular}
\end{center}

The learned top-k-2 contextual adapter clearly beats random fixed top-k-2 and a
norm-matched dense active-rank control. The column ablation fingerprints are
nontrivial:
\[
\operatorname{mean} |\Delta L_{\mathrm{ablate}}| = 0.05213485,
\qquad
\max |\Delta L_{\mathrm{ablate}}| = 0.19186115.
\]
However, rank-matched top-k-1 contextual routing currently has better CE than
the promoted top-k-2 contextual run. Also, the current fingerprint artifact
does not yet measure held-out batch/position fingerprint stability or compare
against top-k-1 fingerprints. The current decision is therefore:
\[
\texttt{causal\_column\_claim\_supported} = \texttt{False}.
\]

\section{Useful Conclusions}

The evidence so far supports the following conclusions.

\begin{enumerate}
  \item The infrastructure is now strong enough to run useful small-scale
  science. Experiments are command-driven, artifact-checked, and have local
  plus Colab evidence for the major promoted changes.

  \item The main positive result is contextual support routing. A contextual
  MLP router using neighboring hidden states and position features is much
  better than a linear router for selecting sparse residual columns. The effect
  is large compared with the earlier \(10^{-4}\)-scale HEP gains.

  \item The router is a major bottleneck. Exhaustive support audits show that a
  per-token oracle support often beats the learned router, and contextual
  selectors can recover a large fraction of this gap.

  \item Support width matters, but it is confounded with active rank and
  residual scale. The top-k-2 gain is real in the current harness, but it is not
  yet proof of cooperative columns.

  \item HEP/settling is currently a stability mechanism, not the source of the
  major gain. Temporal-consistency clipped HEP is safe and repeatable, but its
  improvements are small. In the contextual-router runs, alpha 0 is best.

  \item Causal column structure remains an open question. Column ablation
  fingerprints are nontrivial, learned sparse routing beats random fixed
  support, and norm-matched dense is worse than learned sparse top-k-2.
  Nevertheless, rank-matched top-k-1 contextual CE is currently better, and
  fingerprint stability is not yet established.
\end{enumerate}

\section{What Remains To Be Done}

\subsection{Immediate Deconfounding}

The next immediate step is to extend the causal fingerprint audit:
\begin{itemize}
  \item compute held-out batch and held-out position fingerprint stability;
  \item compute rank-matched top-k-1 contextual fingerprints;
  \item compare top-k-1, top-k-2, random fixed support, and dense controls under
  matched residual norms, active rank, stored parameters, and approximate
  FLOPs;
  \item report both identity support churn and functional churn.
\end{itemize}

\subsection{Better Columnar Evidence}

The project should avoid claiming ``causal columns'' merely from support
utilization or CE. Stronger evidence would include:
\begin{itemize}
  \item stable column intervention fingerprints across batches, positions,
  seeds, and data representations;
  \item low oracle-support regret without collapsing into one dominant pair;
  \item lower functional churn or lower anchor-task drift than matched dense
  adapters;
  \item recurrent reuse of columns under task-free continual learning;
  \item pairwise synergy that remains after rank and norm matching.
\end{itemize}

\subsection{Scaling and Backend}

The Colab path has been useful but brittle. A RunPod bridge is now being set up
so that larger GPU runs can be launched by the same command-driven artifact
contract. The immediate backend goal is not to add new science, but to make
replication and scale-up less fragile.

\subsection{Architectural Variants Worth Testing}

The contextual router result suggests several next architecture tests:
\begin{itemize}
  \item train the contextual support head directly from oracle-support audit
  targets, then fine-tune end-to-end;
  \item compare hard top-k routing to softened or temperature-controlled
  support selection;
  \item test support composers that can exploit pair synergy rather than merely
  selecting columns independently;
  \item add explicit load-balancing only after checking that it does not trade
  off against CE or causal-fingerprint stability;
  \item compare against dense-teacher residual distillation at matched active
  compute.
\end{itemize}

\section{Draft Takeaway for a Colleague}

If another group is experimenting with columnar residual models, the most
portable lesson is this: do not treat sparse-column routing as a minor detail.
In the current RelaLeap harness, a linear hidden-state router left substantial
performance and support-selection headroom on the table. Adding simple
contextual information--previous and next hidden states, local differences, and
position features--made the support router dramatically better. The improvement
replicated across local and Colab backends and across character-level and
tokenized settings. At the same time, the result should not be overinterpreted
as proof that the columns are already clean causal modules. The current best
scientific claim is narrower and stronger:
\[
\boxed{
\text{contextual support selection is a high-leverage mechanism for sparse
columnar residual adapters.}
}
\]
The next burden of proof is to show that this better routing yields stable,
reusable, causally interpretable column functions under matched controls.

\appendix

\section{Primary Artifact References}

\begin{itemize}
  \item \texttt{results/reports/contextual\_support\_router\_decision/decision\_report.md}
  \item \texttt{results/reports/contextual\_support\_router\_promotion\_gate/decision\_report.md}
  \item \texttt{results/reports/contextual\_support\_router\_promotion\_gate\_satisfaction/decision\_report.md}
  \item \texttt{results/reports/support\_width\_deconfounding\_colab\_audit/decision\_report.md}
  \item \texttt{results/reports/validation\_support\_wide\_exhaustive\_support\_audit/decision\_report.md}
  \item \texttt{results/reports/char\_larger\_support\_wide\_seed2\_contextual\_support\_head\_audit/decision\_report.md}
  \item \texttt{results/reports/char\_larger\_support\_wide\_seed3\_contextual\_support\_head\_audit/decision\_report.md}
  \item \texttt{results/reports/post\_promotion\_support\_wide\_promoted\_default/decision\_report.md}
  \item \texttt{results/reports/token\_larger\_causal\_column\_fingerprint\_audit/decision\_report.md}
  \item \texttt{docs/research\_pivot\_2026\_06\_23.md}
\end{itemize}

\end{document}
